\clearpage
\chapter*{ВВЕДЕНИЕ}
\addcontentsline{toc}{chapter}{ВВЕДЕНИЕ}

Расчёт термодинамических свойств воды и водяного пара является базовой задачей в теплоэнергетике, химической технологии и смежных инженерных областях.
Практические расчёты требуют согласованной модели, охватывающей широкую область давлений и температур, а также обеспечивающей воспроизводимые результаты при использовании в различных программных средах.
В качестве промышленного стандарта для таких вычислений широко применяется формулировка IAPWS-IF97~\cite{iapws_if97_2012}, основанная на разбиении области состояния на регионы и использовании специализированных уравнений для каждого из них.

Актуальность работы определяется потребностью в вычислительном инструменте, который:
\begin{itemize}
\item реализует стандарт IAPWS-IF97 в виде единого ядра с безопасным интерфейсом;
\item поддерживает типовые инженерные маршруты расчёта, включая обратные зависимости;
\item может использоваться как в настольных приложениях, так и в сервисной среде;
\item обеспечивает эксплуатационную готовность за счёт воспроизводимой сборки,
тестирования и автоматизированной поставки.
\end{itemize}

Цель работы --- разработать вычислительную модель расчёта термодинамических свойств воды на основе IAPWS-IF97
и реализовать программный комплекс,
предоставляющий результаты расчёта через несколько каналов доставки:
настольное приложение,
настольную оболочку на базе веб-технологий
и сервис.

Для достижения поставленной цели решены следующие задачи:
\begin{itemize}
\item реализовано вычислительное ядро на языке Rust, включающее уравнения регионов IAPWS-IF97, маршрутизацию расчётов, валидацию входных данных и унифицированный формат результата;
\item поддержаны прямые и обратные режимы ввода-вывода по параметрам $p$, $T$, $h$, $s$, $\rho$, $x$;
\item выполнено численное расширение маршрута расчёта по паре $(\rho, T)$ на всей рабочей области за счёт подбора давления методом секущих с последующим расчётом по маршруту $(p, T)$;
\item разработано нативное настольное приложение на FLTK для интерактивного и табличного расчёта, построения диаграмм и экспорта результатов;
\item реализована настольная оболочка на базе Yew, WebAssembly и Tauri
с разделением клиентской и серверной частей
и использованием общего слоя DTO для обмена данными;
\item реализован gRPC-сервис для высокопроизводительных пакетных расчётов с разделением нагрузки ввода-вывода и вычислительной нагрузки, а также с механизмами ограничения нагрузки;
\item подготовлены контейнеризация в Docker и манифесты Kubernetes; настроены средства непрерывной интеграции и доставки, обеспечивающие воспроизводимую сборку и поставку артефактов.
\end{itemize}

Объект исследования --- вода и водяной пар как рабочая среда, описываемая термодинамическими свойствами в области применимости IAPWS-IF97.
Предмет исследования --- вычислительные зависимости и численные методы расчёта термодинамических свойств, а также архитектура программного комплекса с единым ядром и несколькими адаптерами доставки результата.

Методы и средства: формулы и алгоритмы стандарта IAPWS-IF97;
численные методы решения нелинейных уравнений (метод секущих);
язык программирования Rust;
технологии пользовательского интерфейса FLTK, Yew, WebAssembly и Tauri;
gRPC для сервисного контура;
Docker и Kubernetes для контейнеризации и развёртывания;
GitHub Actions для автоматизации непрерывной интеграции и доставки.

Научная новизна и элементы практического вклада работы состоят в построении
архитектуры «одно вычислительное ядро, несколько каналов доставки», а также в
интеграции численного маршрута $(\rho, T)$ в общую схему вычислений, что расширяет
применимость ядра без замены физической основы IAPWS-IF97.

Структура работы включает пять глав.
В первой главе выполнен анализ предметной области и сформулированы требования к программному комплексу.
Во второй главе описаны математическая модель IAPWS-IF97 и используемые численные методы, включая особенности расчётов в различных регионах и в двухфазной области.
В третьей главе рассмотрены архитектура вычислительного ядра на Rust, система типов, маршрутизация и верификация.
В четвёртой главе описаны проектирование кроссплатформенного интерфейса,
реализация FLTK-приложения
и оболочки на базе Yew, WebAssembly и Tauri,
а также визуализация и экспорт данных.
В пятой главе рассмотрены проектирование gRPC-интерфейса, модель параллелизма, контейнеризация, оркестрация и автоматизация непрерывной интеграции и доставки.
В заключении подведены итоги и сформулированы направления дальнейшего развития.
